Suppose that circle $A$ has center $A = (2, 3)$, that circle $B$ has center $B$ and that circle $C$ has center $C = (0, 1)$. Recall that the power of a point  with respect to some circle, $p_A$ is defined to be $r^2_A - d^2_A$, where $r_A$ is the radius of circle $A$, and $d_A$ is the distance from the point to the center of circle $A$. The equations $p_A = ap_B$ and $p_B = bp_C$ graph circles centered at $(2, 4)$ and $(-1,-2)$, respectively for some values of $a$ and $b$. Find the center of the circle with equation $p_A =\dfrac{a}{b}p_C$.
